\rtmtight
\begin{tblr}{width=\textwidth,colspec={|Q[c,m,2.1cm]|X[l,m]|},hlines,vlines,hspan=minimal,rowsep=1.5pt,colsep=3pt,row{1}={bg=headgray,font=\bfseries}}
\SetCell{c,m}\hfil\logo\hfil & \textbf{POLITEKNIK NEGERI MALANG}\par\textbf{JURUSAN TEKNOLOGI INFORMASI}\par\textbf{PROGRAM STUDI: D4 TEKNIK INFORMATIKA} \\
\SetCell[c=2]{l} \textbf{RENCANA TUGAS MAHASISWA (RTM) DAN RUBRIK PENILAIAN} & \\
Nama Mata Kuliah & Pemrograman Berbasis Framework \\
Kode Mata Kuliah & RTI256003 \quad \textbf{sks / Jam perminggu:} 2 SKS/4 jam \quad \textbf{Semester:} 6 \\
Dosen Pengampu & \begin{enumerate}\item Farid Angga Pribadi, S.Kom., M.Kom.\item Dimas Wahyu Wibowo, S.T., M.T.\item Muhammad Unggul Pamenang, S.St., M.T.\item Dian Hanifudin Subhi, S.Kom., M.Kom.\end{enumerate} \\
\end{tblr}
\assessmentblock{Tugas Framework Fundamentals}{Hands-on Practice}{SCPMK0209-04701, SCPMK0603-04702}{Mahasiswa menyelesaikan tugas framework fundamentals sesuai Sub-CPMK dan konteks mata kuliah.}{Menganalisis kebutuhan, mengerjakan artefak praktik/proyek, menguji hasil, mendokumentasikan evidence, dan mempresentasikan keputusan bila diminta.}{Laporan, artefak digital, repository/prototype, evidence pengujian, dan/atau bahan presentasi.}{Pemahaman struktur dan lifecycle framework & 4 & Struktur project, routing, component/controller, lifecycle, dan konfigurasi dasar framework dijelaskan tepat. & Pemahaman struktur cukup, tetapi lifecycle atau konfigurasi belum lengkap. & Tidak memahami struktur framework atau project tidak berjalan. \\ \\
Implementasi fitur dasar framework & 3.5 & Routing, rendering, data binding/form, validasi sederhana, dan reusable component/controller berjalan sesuai tugas. & Fitur dasar berjalan sebagian, tetapi validasi atau komponen belum konsisten. & Fitur dasar tidak berjalan atau tidak mengikuti pola framework. \\ \\
Evidence praktik dan dokumentasi keputusan & 2.5 & Commit, screenshot, konfigurasi, dan catatan keputusan framework terdokumentasi jelas. & Evidence tersedia, tetapi dokumentasi keputusan atau konfigurasi masih terbatas. & Tidak ada evidence praktik atau dokumentasi tidak dapat diverifikasi. \\}

\assessmentblock{UTS Modul Framework Terintegrasi}{UTS praktik}{SCPMK0209-04701, SCPMK0603-04702}{Mahasiswa menyelesaikan uts modul framework terintegrasi sesuai Sub-CPMK dan konteks mata kuliah.}{Menganalisis kebutuhan, mengerjakan artefak praktik/proyek, menguji hasil, mendokumentasikan evidence, dan mempresentasikan keputusan bila diminta.}{Laporan, artefak digital, repository/prototype, evidence pengujian, dan/atau bahan presentasi.}{Kelengkapan modul terintegrasi & 8 & Modul mencakup route, view/component, controller/service, model/data, validasi, dan alur pengguna. & Modul utama berjalan, tetapi validasi, model, atau alur pengguna belum lengkap. & Modul tidak memenuhi fitur integrasi yang diminta. \\ \\
Kualitas arsitektur dan reusable component & 7 & Kode mengikuti pola framework, concern terpisah, komponen reusable, dan konfigurasi mudah dijalankan ulang. & Arsitektur cukup rapi, tetapi masih ada duplikasi atau concern bercampur. & Kode tidak mengikuti pola framework atau sulit dipelihara. \\ \\
Pengujian dan defense modul & 5 & Kasus uji normal/invalid dilakukan, hasil diverifikasi, dan keputusan teknis dapat dipertahankan. & Pengujian terbatas atau defense belum menjelaskan semua keputusan. & Tidak ada pengujian bermakna atau tidak mampu menjelaskan modul. \\}

\assessmentblock{PBL Aplikasi Berbasis Framework}{Project Based Learning}{SCPMK0603-04703}{Mahasiswa menyelesaikan pbl aplikasi berbasis framework sesuai Sub-CPMK dan konteks mata kuliah.}{Menganalisis kebutuhan, mengerjakan artefak praktik/proyek, menguji hasil, mendokumentasikan evidence, dan mempresentasikan keputusan bila diminta.}{Laporan, artefak digital, repository/prototype, evidence pengujian, dan/atau bahan presentasi.}{Kesesuaian solusi dengan kebutuhan proyek & 28 & Kebutuhan pengguna, rancangan data, role, dan prioritas fitur diterjemahkan menjadi scope aplikasi framework. & Scope cukup sesuai, tetapi sebagian kebutuhan, role, atau prioritas fitur belum jelas. & Aplikasi tidak menunjukkan kebutuhan proyek atau scope tidak konsisten. \\ \\
Integrasi fitur dan kualitas implementasi framework & 24.5 & Routing, UI/component, business logic, data access/API, auth, validasi, dan error handling terintegrasi stabil. & Fitur utama terintegrasi, tetapi validasi, error handling, atau struktur data masih kurang. & Fitur utama tidak terintegrasi atau aplikasi sering gagal. \\ \\
Testing, deployment, dan komunikasi hasil & 17.5 & Test/demo data, deployment atau panduan run, repository, evidence, dan presentasi keputusan tersedia lengkap. & Evidence tersedia, tetapi testing/deployment atau presentasi keputusan belum kuat. & Tidak ada evidence memadai atau aplikasi tidak dapat didemonstrasikan. \\}

\begin{tblr}{width=\textwidth,colspec={|Q[l,m,3.2cm]|X[l,m]|},hlines,vlines,hspan=minimal,row{1-3}={bg=headgray},rowsep=1.5pt,colsep=3pt}
Jadwal Pelaksanaan: & Minggu 1--17 sesuai RPS. Setiap evaluasi memiliki RTM dan rubrik multi-kriteria. \\
Lain-Lain Yang Diperlukan: & LMS, repositori/prototype/proyek, perangkat lunak sesuai mata kuliah, dan perangkat presentasi. \\
Pustaka: & Mengacu pustaka utama dan pendukung pada bagian RPS. \\
\end{tblr}
